\documentclass{book}

\usepackage{amsmath}
\usepackage{amsfonts}
\usepackage{amssymb}
\usepackage{amscd}
\usepackage{amsthm}
\usepackage{graphics}
\usepackage{graphicx}
\usepackage{tikz}
\usepackage{tikz-cd}
\usepackage{bm}

\newtheorem{theorem}{Theorem}
\newtheorem{definition}{Definition}
\newcommand{\bracenom}{\genfrac{\lbrace}{\rbrace}{0pt}{}}

\graphicspath{ {./} }

\title{Putnam - Solutions}
\author{Mingruifu Lin}
\date{August 2026}

\begin{document}

\maketitle

\tableofcontents

\chapter{Putnam 1938}

\section*{A1}
Let $A(k) = ak^3 + bk^2 + ck + d$, where $d = M$ since $A(0) = M$. Let us compute the volume by integrating $A(k)$:
$$V = \int_{-\frac{h}{2}}^{\frac{h}{2}} A(k) dk$$
$$= \left[ \frac{a}{4}k^4 + \frac{b}{3}k^3 + \frac{c}{2}k^2 + Mk \right]_{-\frac{h}{2}}^{\frac{h}{2}}$$
$$= \frac{bh^3}{12} + Mh$$
Our goal is to get
$$V = \frac{h(B + 4M + T)}{6}$$
Notice that
$$B = A \left( -\frac{h}{2} \right) = -a \frac{h^3}{8} + b \frac{h^2}{4} - c \frac{h}{2} + M$$
$$T = A \left( \frac{h}{2} \right) = a \frac{h^3}{8} + b \frac{h^2}{4} + c \frac{h}{2} + M$$
thus
$$B + T = \frac{bh^2}{2} + 2M$$
Substituting into our goal, we simplify
$$V = \frac{h(B + 4M + T)}{6} = \frac{h(\frac{bh^2}{2} + 6M)}{6} = \frac{bh^3}{12} + Mh$$
Hence, the volume formula is indeed correct.

We can use this formula to derive the formula for a cone and sphere. For a cone with base radius $r$ and height $h$, the area formula is
$$A(k) = \pi \left( \frac{r}{h} \left( k + \frac{h}{2} \right) \right)^2$$
thus a 2nd degree polynomial, which satisfies the requirements for using the special formula, hence
$$V = \frac{h(B + 4M + T)}{6}$$
$$= \frac{h(\pi r^2 + 4 \pi \left( \frac{r}{2} \right)^2 + 0)}{6}$$
$$= \frac{h(2 \pi r^2)}{6} = \frac{\pi r^2 h}{3}$$
For a sphere with radius $r = \frac{h}{2}$, the cross-sectional area is 
$$A(k) = \pi \left( \frac{h^2}{4} - k^2 \right)$$
thus again a 2nd degree polynomial in $k$, so we apply the special formula to get
$$V = \frac{h(B + 4M + T)}{6}$$
$$= \frac{h(0 + 4 \pi \frac{h^2}{4} + 0)}{6} = \frac{\pi h^3}{6}$$
Expressed in radius $2r = h$, we get
$$V = \frac{4}{3} \pi r^3$$

\section*{A2}
$$V = 2 \cdot \frac{\pi r^2 h}{3} + \pi r^2 h = \frac{5}{3} \pi r^2 h$$
$$A = 2\pi rh + 2 \pi r \sqrt{r^2 + h^2}$$
We want to express either $r$ or $h$ in terms of $A$ and the other:
$$A - 2\pi rh = 2\pi r \sqrt{r^2 + h^2}$$
$$A^2 - 4A\pi rh + 4\pi^2 r^2 h^2 = 4\pi^2 r^2 (r^2 + h^2)$$
$$A^2 - 4A\pi rh + 4\pi^2 r^2 h^2 = 4 \pi^2 r^4 + 4\pi^2 r^2 h^2$$
$$A^2 - 4A\pi rh = 4 \pi^2 r^4$$
$$h = \frac{A^2 - 4\pi^2 r^4}{4A\pi r}$$
Substituting into the volume equation, we get
$$V = \frac{5}{3} \pi r^2 \frac{A^2 - 4\pi^2 r^4}{4A\pi r}$$
$$= \frac{5}{12} \frac{r(A^2 - 4\pi^2 r^4)}{A}$$
We want to maximize this equation given a fixed $A$, so we can ignore the coefficient and denominator, leaving only
$$r(A^2 - 4\pi^2 r^4) = A^2 r - 4\pi^2 r^5$$
Setting the derivative to zero, we get
$$A^2 - 20\pi^2 r^4 = 0$$
$$\Rightarrow r = \sqrt{\frac{A}{\sqrt{20}\pi}}$$
Note that, from the height equation, we know that
$$0 < A^2 - 4\pi^2 r^4$$
$$\Rightarrow r < \sqrt{\frac{A}{2\pi}}$$
$\sqrt{20} > \sqrt{16} = 4 > 2$, so the $r$ we found is safely within the bounds.

\section*{A3}
Velocities are
$$\frac{dx}{dt} = 3t^2 - 1$$
$$\frac{dy}{dt} = 4t^3 + 1$$
I think they are asking for the speed, which is
$$\sqrt{\left( \frac{dx}{dt} \right)^2 + \left( \frac{dy}{dt} \right)^2} = \sqrt{16t^6 + 9t^4 + 8t^3 - 6t^2 + 2}$$
Since we are finding the maximum, we can ignore the square root. Since a polynomial is differentiable everywhere, its maxima happen when derivative is zero
$$96t^5 + 36t^3 + 24t^2 - 12t = 0$$
Clearly, $t = 0$ is an extrema (or saddle). To know whether it's a maximum, we can take the second derivative at $0$
$$480t^4 + 108t^2 + 48t - 12$$
which equals $-12$ (negative) when $t$ are replaced with $0$, hence $t = 0$ is a local maximum.

To know whether it's an inflection point, we can compute $\frac{d^2y}{dx^2}$ at $t = 0$, which should be possible since the tangent is not vertical (specifically, $\vec{v}(0) = (-1, 1)$):
$$\frac{dy}{dx} = \frac{\frac{dy}{dt}}{\frac{dx}{dt}} = \frac{3t^2 - 1}{4t^3 + 1}$$
$$\frac{d}{dx} \frac{dy}{dx} = \frac{d}{dt} \frac{dy}{dx} \cdot \frac{dt}{dx} = \frac{d}{dt} \frac{dy}{dx} \cdot \frac{1}{\frac{dx}{dt}}$$
$$= \frac{(4t^3 + 1)(6t) - (3t^2 - 1)(12t^2)}{(4t^3 + 1)^2} \cdot \frac{1}{3t^2 - 1}$$
$$= \frac{-12t^4 + 12t^2 + 6t}{(4t^3 + 1)^2(3t^2 - 1)}$$
Substituting $0$, the whole thing becomes $0$. Hence, the second derivative is zero, but we still need to know whether there is a sign change for the second derivative (which represents change in concavity).

For $t$ near $0$, the denominator is always negative. However, the derivative of the numerator is $6$ (positive) at $t = 0$, and since the value of the numerator is $0$ at $t = 0$, the numerator changes sign when $t < 0$ becomes $t > 0$. Hence, there is a sign change in the second derivative at $t = 0$, hence it's an inflection point.

\end{document}
